% !TeX root = ../drafts/e-recursion.tex
\section{Recursion: the verifier as a circuit}\label{annex:rec}

A recursion proof shows that leanVM proofs verify. It is not a run of the RISC-V machine: a second machine, the \emph{recursion machine}, proves a fixed list of rows that check the proofs, with the same bus (\S\ref{sec:gp}, \S\ref{sec:gkr}), table sumcheck (\S\ref{sec:air}), Flock (Annex~\ref{annex:flock}) and stacked opening (\S\ref{sec:stacking}, Annex~\ref{annex:pcs}). The rows run the verifier's core, the protocol of \S\ref{sec:e2e-unrolled} short of the evaluations only the program or the circuits fix; those, and a few others the outer verifier is better placed to make, leave the circuit as claims in the outer statement, which the outer verifier settles after the outer proof (\S\ref{rec:statement}).

\subsection{Shapes and circuits}\label{rec:shape}

A \emph{shape} is a program $P$ with its digest, the eleven table heights $\tau_j$ and the rate. It fixes every loop bound of the core: the layout of every stack, the bus's depth $\mu$, the rounds of every sumcheck, the WHIR ladder, its query counts and the depth of every Merkle path. A \emph{circuit} $\rcirc$ for a list of shapes is the list of rows the core executes on one proof of each, and which of their slots carry the same value. It is a function of the shapes alone: the verifier builds it from them, never from a proof, by the same code the prover runs, with every value read from the proof replaced by zero. What depends on a proof is only the prover's assignment of values. In particular no row exists or is skipped according to a value: a query index selects nothing, it is a list of bits that the Merkle rows take as the direction of each level, and the paths are hashed in full rather than pruned.

\subsection{Wires and copy cycles}\label{rec:wires}

A \emph{wire} holds one of three kinds of value: a word of $\K$, an element of $\E$ as its three limbs, or a digest as four words. Every row of every table names one wire per slot, and a slot $s$ of row $z$ carries its wire's value $v$ as the tuple $\tup{\slotkey,v_0,v_1,v_2,v_3}$, unused limbs zero, where $\slotkey(s,z)$ is the integer $\intk{2^{32}g+z}$, $g$ the slot's index among all tables' slots: an index column (\S\ref{sec:idxcol}). The slots naming one wire form a cycle in a fixed order, and the slot \pull{}s its tuple at its own key and \push{}es it at $\slotnext(s,z)$, the key of the next slot of the cycle, a public column of the circuit.

\begin{lemma}\label{lem:rec-cycles}
If the bus balances, every slot of a cycle carries the same value.
\end{lemma}
\begin{proof}
Every key is pulled once and pushed once, by distinct slots, so the multiset equality matches the tuple pushed at a key with the one pulled there: the value slot $s$ carries is the value its successor carries. Around a cycle they are all one.
\end{proof}

A wire on one slot is a cycle of one, whose tuple cancels itself: a free value, which is how a scalar the prover sends enters, and how a hint does. An equality the verifier checks merges two wires' cycles into one. A constant, and a word of the statement, is a row of the \textsf{PUB} table, a framework block (\S\ref{sec:leafstack}) whose every coordinate is public: the verifier knows its value. Padding rows have $\slotnext=\slotkey$ in every slot, so whatever they carry cancels, and satisfy their table's identities with zero columns or, for a hash row, as an honest compression.

\subsection{Tables}\label{rec:tables}

Each table's slots are degree-two forms over its columns (\S\ref{sec:m3}), so an output need not be committed: it rides the bus as the form that computes it.
\begin{itemize}[itemsep=1pt,topsep=3pt]
\item \textsf{EMUL}: $c=a\cdot b+d$ over $\E$. Columns $a,b,d$, nine words of $\K$. The slot of $c$ carries $c_0=a_0b_0+a_1b_2+a_2b_1+d_0$, $c_1=a_0b_1+a_1b_0+a_1b_2+a_2b_1+a_2b_2+d_1$, $c_2=a_0b_2+a_1b_1+a_2b_0+a_2b_2+d_2$, the product modulo $y^3+y+1$.
\item \textsf{EXK}: $c=a\cdot k+d$ with $k\in\K$. Seven columns; $c_i=a_ik+d_i$.
\item \textsf{HASH}: one BLAKE2s compression, its words ports of a packed witness of the \textsf{HASH} class circuit (\S\ref{flock:classes}), which Flock proves, plus a selector bit $b$, with the identity $b^2=b$. Its slots are the chaining value $h$, the counter and the finalization word, the message's first half through the selector ($m_k+b(m_k+m_{4+k})$, the left half when $b=0$ and the right one when $b=1$), $b$, the message words $m_4,m_5,m_6$ as an element of $\E$, $m_7$, every message word alone, the output, and the output's first three words as an element of $\E$. A transcript step wires $h$ to the parameter IV and the counter to $(64,\text{final})$; a Merkle node too, its child on the selector slot and its sibling free; a block of a leaf's hash wires $h$ to the previous block's output and its counter to the block's. The hash rows are split, in order, over several tables of this one shape, the first ones taking the largest powers of two that fit, so that padding to a power of two wastes little.
\item \textsf{SPLIT}: a word and its 64 bits, with $w=\sum_i\intk{2^i}\,b_i$ and $b_i^2=b_i$. Read one way it splits a word, read the other it packs bits into one.
\item \textsf{CAST}: four words seen as a digest, as an element of $\E$, as two elements of $\E$ with a zero top limb (a digest's two halves on the stream), and as four words.
\end{itemize}
Each table's summand in the table sumcheck is its two bus forms, weighted $1$ and $\xi$ as the RISC-V tables', and its identities, each with a power of $\xi$ of its own from $\xi^2$, the tables' ranges disjoint, so that the target pins each identity's sum to zero.

\subsection{The core in rows}\label{rec:core}

The Fiat-Shamir chain is a chain of \textsf{HASH} rows: absorbing a scalar is a compression of the chaining value with the scalar and the observe tag (\S\ref{sec:e2e-fiat-shamir}), and a challenge is the first three words of a compression with the squeeze tag. The scalars hashed are the wires the arithmetic uses, so what is hashed is what is used. A round polynomial's fixed coefficient is derived from the claim, the others read and absorbed in order; a proof of work is two compressions and a split of the digest's first word, its low bits held to zero. The announced heights and rate are held to the shape's; the final clock is split and held to a live clock at slot zero. The sumchecks, the bus's GKR, its decomposition and the table sumcheck's final identity are \textsf{EMUL} and \textsf{EXK} rows; the tables' bus forms are evaluated at the claimed column values, a product coordinate as a product of two claimed values. Each Flock reduction is replayed up to its matrices. The opening's query indices are the bits of challenges, a query's bits possibly straddling two limbs; each query's leaf is hashed from the shared zero prefix of the absent lanes (\S\ref{sec:stacking}) and its path in full, the direction of each level the query's bit there, and the root is merged with the one on the stream. The query's index as a word, for the induced weights, is packed from its bits.

\subsection{The outer statement}\label{rec:statement}

The outer statement is, for each inner proof, its output, its shape, and what the core leaves:
\begin{itemize}[itemsep=1pt,topsep=3pt]
\item the program claim: the table sumcheck's final identity short of the bytecode producer's program columns and of RAM's image, as one evaluation of the stacked bytecode table at the sumcheck's point and the fingerprint's, twisted by the multiplicity bits' Frobenius powers, and one of the image;
\item per Flock circuit, its matrix form $\erow^{\top}(A_0+\lcalpha B_0)\wcol$ with its points and slices;
\item the ring-switched claims' shares of the opening's target and of its terminal weight (\S\ref{sec:ringswitch}), with the map's and the batching challenges, the terminal point, each Flock claim's outer coordinates and the multiplicity bits, which with the matrix forms locate every ring-switched claim.
\end{itemize}
Every word of it is a \textsf{PUB} row on the wire the core computed. The outer transcript's initial state is a digest of the shapes, its first block a digest of the statement. The outer verifier rebuilds $\rcirc$ from the shapes, checks the outer proof against the statement, then evaluates each program claim on the program, each matrix form on its circuit, and each ring-switched share from the claims by the native closed forms.

\paragraph{Soundness.} If the outer proof verifies then, except with the probability its own soundness allows, the bus balances, every identity holds and every hash row is a compression (Annex~\ref{annex:flock}). By Lemma~\ref{lem:rec-cycles} the values form one consistent assignment of the circuit's wires. The free wires on the stream's scalars and the Merkle hints are then an inner proof on which every check of the core holds, with challenges that are the transcript's compressions of exactly those scalars, leaving exactly the claims the statement holds; the outer verifier's settlement of those claims is the rest of the inner verifier. So each inner proof verifies, against the output and the program the statement names. As for every recursive composition of Fiat-Shamir proofs, knowledge soundness of the composition rests on the transcript hash behaving as a random oracle.
